% =====================================================================
%  M3 -- Step-size stability on a coupled quadratic
% =====================================================================
\subsection{M3: Step-size stability on a coupled quadratic}
\setprob{M3}

\begin{frame}{\probtitle{M3}{Step-size stability on a coupled quadratic}}
  \begin{probbox}[Problem M3 \hfill Mathematical \quad $\star\star$]
    Let $f(\x) = \tfrac12\x\T\bA\x - \vb\T\x$ with
    $\bA = \begin{pmatrix}3&1\\1&3\end{pmatrix}$, $\vb = \begin{pmatrix}4\\4\end{pmatrix}$,
    and run GD, $\x_{k+1} = \x_k - \alpha\nabla f(\x_k)$.
    \begin{enumerate}[(a)]
      \item Find $\x^\star$ and show that $\e_k = \x_k - \x^\star$ obeys $\e_{k+1} = (\bI - \alpha\bA)\,\e_k$.
      \item Use the eigenpairs of $\bA$ to find \emph{all} $\alpha$ for which GD converges from every $\x_0$.
      \item Find the $\alpha$ that makes the worst-case contraction factor smallest, and that factor.
      \item With this $\alpha$ and $\x_0 = (3,0)$, compute $\x_1,\x_2$ and check that $\norm{\e_2} = \norm{\e_0}/9$.
      \item How many iterations guarantee $\norm{\e_k}\le 10^{-6}\norm{\e_0}$?
      \item Describe exactly what happens when $\alpha = 0.6$.
    \end{enumerate}
  \end{probbox}
  \footnotesize\textcolor{mGrey}{What this practises: an exact stability analysis, the best fixed step, and comparing step choices.}
\end{frame}

\begin{frame}{\soltitle{M3}{1}{4}}
  \textbf{(a)} Solve $\nabla f(\x) = \bA\x - \vb = \mathbf 0$. Since $\det\bA = 8$,
  \[
    \x^\star = \bA^{-1}\vb = \frac18\begin{pmatrix}3&-1\\-1&3\end{pmatrix}\begin{pmatrix}4\\4\end{pmatrix}
    = \ans{\begin{pmatrix}1\\1\end{pmatrix}} .
  \]
  Using $\vb = \bA\x^\star$:
  $\;\e_{k+1} = \x_k - \alpha(\bA\x_k - \bA\x^\star) - \x^\star = (\bI - \alpha\bA)\,\e_k .$

  \vspace{4pt}
  \textbf{Eigen-decomposition.} $\det(\bA-\lambda\bI) = (3-\lambda)^2 - 1 = 0$ gives
  \[
    \lambda_1 = 2,\;\; \vv_1 = \tfrac1{\sqrt2}(1,-1)\T; \qquad
    \lambda_2 = 4,\;\; \vv_2 = \tfrac1{\sqrt2}(1,1)\T .
  \]
  Writing $\e_k = c_1^{(k)}\vv_1 + c_2^{(k)}\vv_2$ splits the iteration into two separate scalar ones:
  \[
    c_i^{(k)} = (1 - \alpha\lambda_i)^k\, c_i^{(0)}, \qquad i = 1,2 .
  \]
\end{frame}

\begin{frame}{\soltitle{M3}{2}{4}}
  \begin{columns}[T]
    \begin{column}{0.5\textwidth}
      \textbf{(b)} To converge for \emph{every} $\e_0$ we need $|1-\alpha\lambda_i|<1$ for both $i$:
      \[
        0<\alpha<1 \;\text{ and }\; 0<\alpha<\tfrac12
      \]
      \[
        \Longrightarrow\;\ans{0 < \alpha < \tfrac{2}{\lmax} = \tfrac12}
      \]
      \textbf{(c)} $\rho(\alpha) = \max\{|1-2\alpha|,\,|1-4\alpha|\}$ is smallest where the two branches cross:
      \[
        1 - 2\alpha = -(1-4\alpha) \Rightarrow \ans{\alpha^\star = \tfrac13},\;\; \ans{\rho^\star = \tfrac13}
      \]
      \small In general $\alpha^\star = \frac{2}{\lmin+\lmax}$ and $\rho^\star = \frac{\kappa-1}{\kappa+1}$; here $\kappa = 2$.
    \end{column}
    \begin{column}{0.48\textwidth}
      \centering
      \includegraphics[width=\linewidth]{fig_amplification}
    \end{column}
  \end{columns}
\end{frame}

\begin{frame}{\soltitle{M3}{3}{4}}
  \textbf{(d)} With $\alpha = \tfrac13$ and $\x_0 = (3,0)$:
  \begin{align*}
    \g_0 &= \bA\x_0 - \vb = (9,3) - (4,4) = (5,-1), &
    \x_1 &= (3,0) - \tfrac13(5,-1) = \ans{\bigl(\tfrac43,\tfrac13\bigr)}\\
    \g_1 &= \bigl(4+\tfrac13,\ \tfrac43+1\bigr) - (4,4) = \bigl(\tfrac13,-\tfrac53\bigr), &
    \x_2 &= \bigl(\tfrac43-\tfrac19,\ \tfrac13+\tfrac59\bigr) = \ans{\bigl(\tfrac{11}9,\tfrac89\bigr)}
  \end{align*}
  So $\e_0 = (2,-1)$ and $\e_2 = \bigl(\tfrac29,-\tfrac19\bigr) = \tfrac19\,\e_0$, which gives $\norm{\e_2} = \norm{\e_0}/9$. \checkmark

  \vspace{4pt}
  \textbf{Why exactly $1/9$?} Split $\e_0 = \tfrac12(1,1) + \tfrac32(1,-1)$. The factors are
  $1-\tfrac43 = -\tfrac13$ on $(1,1)$ and $1-\tfrac23 = \tfrac13$ on $(1,-1)$. Both have size $\tfrac13$,
  so \emph{every} component shrinks by $\tfrac13$ per step: the optimal step \alert{balances} the two modes.
  \[
    \e_1 = -\tfrac16(1,1) + \tfrac12(1,-1) = \bigl(\tfrac13,-\tfrac23\bigr) \;\checkmark
  \]
\end{frame}

\begin{frame}{\soltitle{M3}{4}{4}}
  \textbf{(e)} $\norm{\e_k} = 3^{-k}\norm{\e_0}$, so we need $3^{-k} \le 10^{-6}$:
  \vspace{-4pt}
  \[
    k \ge \frac{6\ln 10}{\ln 3} = \frac{13.8155}{1.0986} = 12.58
    \;\Longrightarrow\; \ans{k = 13}
  \]
  \textbf{(f) $\alpha = 0.6 > 0.5$.} The factors are $1 - 0.6\cdot2 = -0.2$ and $1 - 0.6\cdot4 = \alert{-1.4}$:
  \[
    \e_k = \tfrac12(-1.4)^k(1,1) + \tfrac32(-0.2)^k(1,-1).
  \]
  \begin{itemize}\small
    \item The $(1,1)$ mode \alert{grows} $\times1.4$ per step, flipping sign each time: $\e_6 \approx (3.765, 3.765)$.
    \item The $(1,-1)$ mode still decays. One stiff direction ($\lmax$) sets the step limit.
  \end{itemize}
  \begin{keybox}[Lesson]
    The blow-up happens along the eigenvector of $\lmax$, just as it does for explicit Euler on a stiff ODE.
  \end{keybox}
\end{frame}
